% =====================================================================
\chapter*{Parte B. Problemas}
\phantomsection\addcontentsline{toc}{chapter}{Parte B. Problemas}
\markboth{Parte B. Problemas}{}
\label{cap:parteB}
% =====================================================================

Seis problemas de desarrollo (70 puntos en total). La puntuación de cada apartado figura entre paréntesis. Las soluciones están en el \cref{sec:sol-problemas}; el título de cada problema enlaza con la suya.

% ---------------------------------------------------------------------
\section*{Problema 1. Retropropagación en un perceptrón de dos capas (18 puntos)}
\phantomsection\addcontentsline{toc}{section}{Problema 1. Retropropagación en un perceptrón de dos capas}

\begin{ejercicio}{3}{eje:mlp}
\textbf{Módulos M08 y M09.} Una red $2\to 2\to 3$ con activación ReLU procesa un lote de dos ejemplos. La pérdida es la entropía cruzada media del lote:
\begin{gather*}
Z^{(1)}=XW_1+b_1,\qquad H=\ReLU\bigl(Z^{(1)}\bigr),\qquad Z^{(2)}=HW_2+b_2,\\
P=\softmax\bigl(Z^{(2)}\bigr)\ \text{por filas},\qquad L=\frac12\sum_{i=1}^{2}\bigl(-\ln P_{i,y_i}\bigr),
\end{gather*}
con
\begin{gather*}
X=\begin{pmatrix}2&1\\1&-1\end{pmatrix},\qquad y=(0,\,2),\qquad
W_1=\begin{pmatrix}1&0{,}5\\-1&1\end{pmatrix},\qquad b_1=(0;\ -0{,}5),\\
W_2=\begin{pmatrix}1&0&-1\\0&1&1\end{pmatrix},\qquad b_2=(0;\ 0;\ 0).
\end{gather*}
El primer ejemplo es de la clase~0 y el segundo, de la clase~2. Sea $Y$ la matriz de etiquetas en codificación \emph{one-hot}.
\begin{apartados}
  \item \pts{3} Realice la pasada hacia delante: calcule $Z^{(1)}$, $H$, $Z^{(2)}$, $P$, la pérdida $\ell_i$ de cada ejemplo y $L$.
  \item \pts{3} Para el primer ejemplo, calcule el jacobiano de la softmax, $J=\diag(p)-p\T p$, y el vector fila $\partial\ell_1/\partial p$. Compruebe que su producto es $p-e_{y_1}$ y demuestre en general que, si $\hat y=\softmax(\theta)$ y $\mathrm{CE}(y,\hat y)=-\sum_k y_k\ln\hat y_k$ con $y$ \emph{one-hot}, entonces $\partial\mathrm{CE}/\partial\theta=\hat y-y$. Explique de dónde procede el factor $1/2$ de $\Delta^{(2)}=(P-Y)/2$.
  \item \pts{3} Calcule $\Delta^{(2)}$, $\nabla W_2=H\T\Delta^{(2)}$ y $\nabla b_2$ (suma de las filas de $\Delta^{(2)}$). Compruebe las formas.
  \item \pts{3} Retropropague: $\nabla H=\Delta^{(2)}W_2\T$, $\Delta^{(1)}=\nabla H\odot\uno\bigl[Z^{(1)}>0\bigr]$, $\nabla W_1=X\T\Delta^{(1)}$ y $\nabla b_1$. Identifique la única de las cuatro unidades ocultas (dos neuronas por dos ejemplos) que está apagada y demuestre que la segunda columna de $\nabla W_1$ solo recibe la contribución del primer ejemplo.
  \item \pts{2} Obtenga la expresión de $\nabla_X L$ en función de $\Delta^{(1)}$ y $W_1$, y calcúlela. ¿Qué componente de la entrada tiene más influencia sobre la pérdida?
  \item \pts{2} Dé un paso de descenso de gradiente con $\eta=0{,}5$ sobre los cuatro tensores de parámetros y recalcule la pérdida de cada ejemplo y la media. ¿Mejoran los dos ejemplos?
  \item \pts{2} Para una red con una capa oculta de este tipo, con entrada de dimensión $D_x$, $H_o$ unidades ocultas y salida de dimensión $D_y$, cuente los parámetros y aplique la fórmula a esta red. Explique qué cambia en las fórmulas anteriores si la pérdida es la \emph{suma} del lote en lugar de la media, y por qué se prefiere la media cuando se cambia el tamaño del lote.
\end{apartados}
\fuente{adaptado de Stanford CS224N, invierno de 2018, tarea~1, problema~2 «Neural Network Basics», apartados (b), (c) y~(d)~\cite{cs224n2018}. Se sustituye la sigmoide por ReLU, se añade un lote de dos ejemplos y los datos numéricos son propios.}
\end{ejercicio}

% ---------------------------------------------------------------------
\section*{Problema 2. Optimización (10 puntos)}
\phantomsection\addcontentsline{toc}{section}{Problema 2. Optimización}

\begin{ejercicio}{2}{eje:opt}
\textbf{Módulo M07.} En este problema los vectores son columnas.
\begin{apartados}
  \item \pts{5} Sea $J(w)=w\T Aw+b\T w$ con
  \[
  A=\begin{pmatrix}2&0\\0&1/2\end{pmatrix},\qquad b=\begin{pmatrix}-4\\-1\end{pmatrix},
  \]
  y sea el descenso de gradiente $w^{(k+1)}=w^{(k)}-\eta\nabla J\bigl(w^{(k)}\bigr)$.
  (i)~Calcule $\nabla J$ y el minimizador $w^{*}$.
  (ii)~Demuestre que el error $e^{(k)}=w^{(k)}-w^{*}$ cumple $e^{(k+1)}=(I-2\eta A)\,e^{(k)}$.
  (iii)~Deduzca la condición sobre $\eta$ para que el método converja desde cualquier punto de partida y calcule el número de condición $\kappa$ de la hessiana.
  (iv)~Calcule la tasa $\eta^{*}=2/(\mu+L)$ y el factor de contracción que produce.
  (v)~Con $\eta^{*}$, dé dos pasos desde $w^{(0)}=(0,0)\T$.
  \item \pts{3} AdamW con un único parámetro: $\theta_0=2$, $\eta=0{,}1$, $\lambda=0{,}1$, $\beta_1=0{,}9$, $\beta_2=0{,}999$ y $\varepsilon=0$. Los gradientes son $g_1=4$ y $g_2=1$. Calcule $m_t$, $v_t$, $\hat m_t$, $\hat v_t$ y $\theta_t$ para $t=1$ y $t=2$.
  \item \pts{2} Un calendario combina calentamiento lineal durante $W=500$ pasos y decaimiento coseno hasta $T=5500$, con $\eta_{\max}=10^{-3}$ y $\eta_{\min}=10^{-4}$. Calcule $\eta_t$ en $t=250$, $3000$, $4250$ y $5500$.
\end{apartados}
\fuente{el apartado~(a) está adaptado de Berkeley CS189, primavera de 2023, examen parcial, pregunta~2 «Convergence of Gradient Descent»~\cite{cs189mid23}; se particulariza a una matriz diagonal y se añaden las preguntas (iii)--(v). Los apartados (b) y~(c) son de elaboración propia.}
\end{ejercicio}

% ---------------------------------------------------------------------
\section*{Problema 3. Aprendizaje supervisado (12 puntos)}
\phantomsection\addcontentsline{toc}{section}{Problema 3. Aprendizaje supervisado}

\begin{ejercicio}{2}{eje:sup}
\textbf{Módulo M08.}
\begin{apartados}
  \item \pts{4} Ajuste por mínimos cuadrados la recta $y=\theta_0+\theta_1x$ a los datos $x=(1,2,3)$, $y=(1,2,6)$. Calcule $X\T X$, $X\T y$, $\hat\theta$, los residuos, la suma de cuadrados residual $\SSE$ y la varianza de máxima verosimilitud $\hat\sigma^{2}$ del ruido gaussiano.
  \item \pts{4} \emph{Regresión logística con un único parámetro.} Cada ejemplo tiene una característica $x\in\R$ y una etiqueta $y\in\{0,1\}$. El modelo es $P(y=1\mid x)=\sigma(x-\alpha)$, de modo que la frontera de decisión es $x=\alpha$. Para un ejemplo, la pérdida es
  \[
  J(\alpha)=-y\ln\sigma(x-\alpha)-(1-y)\ln\bigl(1-\sigma(x-\alpha)\bigr).
  \]
  (i)~Demuestre que $J'(\alpha)=y-\sigma(x-\alpha)$ y escriba la regla del SGD con paso $\epsilon>0$. Aplíquela una vez con $\epsilon=0{,}5$, desde $\alpha=1$, al ejemplo $(x,y)=(3,1)$.
  (ii)~Calcule $J''(\alpha)$ y decida si $J$ es convexa en $\R$.
  (iii)~Para los dos ejemplos $(0,0)$ y $(2,1)$, demuestre que la suma de las pérdidas se minimiza en $\alpha^{*}=1$.
  (iv)~Dé un conjunto de datos para el que la suma de las pérdidas no tenga mínimo en $\R$ y justifíquelo.
  \item \pts{2} Una regresión softmax devuelve los logits $z=(1,\,0,\,-1,\,2)$ y la clase correcta es la~3. Calcule $p$, la pérdida y el gradiente respecto a $z$.
  \item \pts{2} Con $1200$ ejemplos se comparan 5 valores de $\lambda$ mediante validación cruzada con $k=4$ particiones. ¿Cuántos ejemplos tiene cada partición? ¿Con cuántos se entrena en cada ronda? ¿Cuántos entrenamientos se hacen en total, contando el reentrenamiento final? ¿Por qué no se elige $\lambda$ con el conjunto de prueba?
\end{apartados}
\fuente{el apartado~(b) está adaptado de Berkeley CS189, primavera de 2019, examen parcial, pregunta~5 «Logistic Regression with One Feature»~\cite{cs189mid19}; los datos numéricos de (i) y~(iii) son propios. Los apartados (a), (c) y~(d) son de elaboración propia.}
\end{ejercicio}

% ---------------------------------------------------------------------
\section*{Problema 4. Diferenciación automática (10 puntos)}
\phantomsection\addcontentsline{toc}{section}{Problema 4. Diferenciación automática}

\begin{ejercicio}{2}{eje:ad}
\textbf{Módulo M10.} Con entradas $a$, $b$ y $c$, se evalúa la expresión
\[
d=ab,\quad e=\sigma(d),\quad f=\ReLU(a),\quad g=3a,\quad h=2e+f+g,\quad i=ch,\quad j=f+i^{2},
\]
en el punto $a=1$, $b=0$, $c=0{,}5$. Se usa $\ReLU'(t)=\uno[t\ge 0]$.
\begin{apartados}
  \item \pts{2} Dibuje el grafo computacional, indique qué nodos tienen más de una salida y realice la pasada hacia delante.
  \item \pts{3} Aplique el modo inverso con adjuntos, empezando por $\bar\jmath=1$, hasta obtener $\partial j/\partial a$, $\partial j/\partial b$ y $\partial j/\partial c$. Escriba también $\partial j/\partial a$ como expresión simbólica en las variables $a,\dots,i$, sin símbolos de derivada.
  \item \pts{2} Repita el cálculo de $\partial j/\partial a$ en modo directo con el número dual $a=1+\varepsilon$, usando $\sigma(u+v\varepsilon)=\sigma(u)+\sigma'(u)\,v\varepsilon$ y $\ReLU(u+v\varepsilon)=\ReLU(u)+\uno[u\ge 0]\,v\varepsilon$. ¿Cuántas pasadas necesitaría cada modo para obtener el gradiente completo respecto a $(a,b,c)$?
  \item \pts{3} El modelo \texttt{MLPCaracteres} de M10, §6.2, tiene contexto $B=3$, embeddings de dimensión $d=10$, $H=64$ unidades ocultas y $\abs{\V}=28$, con $N=4084$ parámetros. Con un lote de 128 ejemplos, el modo inverso guarda la matriz de embeddings concatenados $E$ ($128\times 30$), $z^{(1)}$, $h$, $z^{(2)}$ y $p$. Cuente cuántos números ocupan y compárelo con $N$. Para la capa oculta, calcule los FLOP de la pasada hacia delante y de la pasada hacia atrás (2~FLOP por multiplicación-suma) y relaciónelos con la regla $C\approx 6ND$.
\end{apartados}
\fuente{los apartados (a) y~(b) están adaptados de Berkeley CS189, primavera de 2023, examen final, pregunta~4 «Backpropagation on an Arithmetic Expression»~\cite{cs189fin23}, que pide las mismas derivadas en forma simbólica; aquí se añaden valores numéricos y el modo directo. El apartado~(d) es de elaboración propia.}
\end{ejercicio}

% ---------------------------------------------------------------------
\section*{Problema 5. Entrenamiento de redes profundas (10 puntos)}
\phantomsection\addcontentsline{toc}{section}{Problema 5. Entrenamiento de redes profundas}

\begin{ejercicio}{2}{eje:prof}
\textbf{Módulo M11.}
\begin{apartados}
  \item \pts{3} Demuestre que $\softmax(z+c\uno)=\softmax(z)$ para todo $c\in\R$. Para los logits $z=(90,\,88,\,80)$, explique qué ocurre si la softmax se calcula de forma ingenua en fp32. Calcule $\LSE(z)$ de forma estable, $p$ y la pérdida si la clase correcta es la~1 (logit~88).
  \item \pts{3} Normalice $x=(-1,\,0,\,2,\,3)$ con LayerNorm y con RMSNorm ($\gamma=g=\uno$, $\beta=0$, $\epsilon=0$). Compruebe que $\RMSNorm(10x)=\RMSNorm(x)$. ¿Ocurre lo mismo con LayerNorm? ¿Y al sumar una constante a todas las componentes?
  \item \pts{2} Recorte el gradiente $g=(2,\,-1,\,2)$ a norma global máxima $c=1{,}5$. Una pila de 8 capas $\tanh$ escalares tiene todas sus preactivaciones en $z=2$. Calcule el factor del gradiente sin conexiones residuales, $\prod_{\ell=1}^{8}\tanh'(2)$, y con ellas, $\bigl(1\pm\tanh'(2)\bigr)^{8}$, para los dos signos.
  \item \pts{2} En bf16, entre $0{,}5$ y $1$ los números representables están separados por $2^{-8}$. Se aplica a un peso $w=1$, por separado, cada una de las actualizaciones $-0{,}001$, $-0{,}003$ y $-0{,}005$ (tres casos independientes). ¿Qué valor se obtiene en bf16 en cada caso? ¿Qué remedio usa la precisión mixta? Por otra parte, 16 bloques residuales suman a un flujo de varianza inicial~1 contribuciones incorreladas de varianza~1 cada una. ¿Qué desviación típica tiene el flujo final?
\end{apartados}
\fuente{la demostración del apartado~(a) está adaptada de Stanford CS224N, invierno de 2018, tarea~1, problema~1(a)~\cite{cs224n2018}. El resto del problema es de elaboración propia.}
\end{ejercicio}

% ---------------------------------------------------------------------
\section*{Problema 6. Estimación de parámetros, cómputo y memoria (10 puntos)}
\phantomsection\addcontentsline{toc}{section}{Problema 6. Estimación de parámetros, cómputo y memoria}

\begin{ejercicio}{1}{eje:est}
\textbf{Módulos M07, M09, M10 y M11.}
\begin{apartados}
  \item \pts{3} Se amplía el modelo de lenguaje de caracteres: contexto $B=5$, $\abs{\V}=28$, $d=16$, $H=128$, con sesgos. Cuente sus parámetros por capa. Estime los FLOP de 800 pasos con minibatches de 128 ejemplos mediante $C\approx 6ND$.
  \item \pts{4} Se pasa a nivel de palabra: $\abs{\V}=50\,000$, $d=64$, contexto de 8 palabras y $H=512$. Cuente los parámetros por capa y el total. ¿Qué fracción está en la capa de salida? Estime la memoria de entrenamiento con AdamW en fp32 (pesos, gradientes, $m$ y $v$), sin contar las activaciones.
  \item \pts{3} Se entrena el modelo de~(b) con $D=10^{8}$ tokens en un equipo que alcanza $10^{12}$~FLOP/s efectivos. Estime el tiempo. ¿Cabe un modelo de 7000 millones de parámetros, con la misma receta de 16 bytes por parámetro, en una GPU de 80~GB?
\end{apartados}
\end{ejercicio}
